\chapternoindex{Основные понятия и сокращения}

В работе используются следующие термины и сокращения:

\begin{description}
    \item[API (Application Programming Interface)] – интерфейс программирования приложений.
    \item[B-Rep (Boundary Representation)] – метод представления трёхмерных тел в САПР.
    \item[\texttt{bossExtrusion}] – операция выдавливания (добавление материала).
    \item[CAD (Computer-Aided Design) / САПР] – система автоматизированного проектирования.
    \item[\texttt{chamfer}] – операция создания фаски.
    \item[\texttt{cutExtrusion}] – операция вырезания выдавливанием.
    \item[DGCNN (Dynamic Graph CNN)] – архитектура нейронной сети для облаков точек.
    \item[\texttt{fillet}] – операция скругления рёбер.
    \item[JSON (JavaScript Object Notation)] – текстовый формат обмена данными.
    \item[MeshCNN] – нейросетевая архитектура для полигональных сеток.
    \item[OBJ] – формат файла для представления трёхмерной геометрии.
    \item[PointNet] – архитектура для классификации и сегментации облаков точек.
    \item[STEP (Standard for the Exchange of Product model data)] – стандарт обмена данными между САПР.
    \item[STL (Stereolithography)] – формат для хранения триангулированных моделей.
    \item[Steps.txt] – промежуточный текстовый формат описания операций.
    \item[Гибридная архитектура] – сочетание нейросетевых методов и детерминированных алгоритмов CAD-ядра.
    \item[Детерминированный алгоритм] – алгоритм с однозначным результатом для одних и тех же входных данных.
    \item[История построения] – последовательность формообразующих операций и их параметров.
    \item[КОМПАС-3D] – система трёхмерного моделирования компании АСКОН.
    \item[Обратный инжиниринг] – процесс воссоздания модели по готовому образцу.
    \item[Сегментация 3D-модели] – разбиение поверхности на области, соответствующие операциям.
\end{description}

